\section{Criticality analysis for NTK tensors}
\label{app:criticality_analysis}
In this section, we will show that the tensors $D$, $A$ and $B$, which describe the NTK statistics at leading order in $\frac{1}{n}$, as well as their higher-order gerenalizations, are also set to criticality once the susceptibilities of section \ref{sec:stability} are constrained to satisfy \eqref{criticalityconstraintsequiv}, \eqref{criticalityconditionchicirc}. We develop our analysis in the order mentioned above:

\begin{enumerate}[label=(\roman*)]

    \item As discussed in \ref{sec:stability}, once the NNGP, quartic vertex and NTK are shown to be stable, the stability of the tensor $D$ follows from the variation of the same tensor in the previous layer. In diagrams:  
\begin{eqnarray}
\begin{tikzpicture}[baseline=(d.base)]
\begin{feynman}
\vertex (d) {};
\vertex[left = 4pt of d, dot, minimum size=0pt] (v12) {};
\vertex[left = 16pt of v12, dot, minimum size=0pt] (l) {};
\vertex[below = 16pt of l, dot, minimum size=3pt, label = {below: {\footnotesize $1$}}] (x1) {};
\vertex[above = 16pt of l, dot, minimum size=3pt, label = {above: {\footnotesize $2$}}] (x2) {};
\vertex[right = 4pt of d, dot, minimum size=0pt] (v34) {};
\vertex[right = 16pt of v34, dot, minimum size=0pt] (r) {};
\vertex[above = 16pt of r, color1, dot, minimum size=3pt, label = {above: {\footnotesize $\textcolor{color1}{3}$}}] (x3) {};
\vertex[below = 16pt of r, color1, dot, minimum size=3pt, label = {below: {\footnotesize $\textcolor{color1}{4}$}}] (x4) {};
\diagram*{
	(x1) -- (v12) -- (x2), 
	 (x3) -- [color1, ghost] (v34) -- [color1, ghost] (x4)
};
\tikzfeynmanset{every blob = {/tikz/fill=black!50, /tikz/minimum size=15pt}}
\vertex[right = 0pt of d, blob, minimum size = 8pt] (dd) {{\scriptsize $\delta$}};
\end{feynman}
\end{tikzpicture}
&=& \sum_{j_1, j_2}
%
\begin{tikzpicture}[baseline=(b)]
\tikzfeynmanset{every blob = {/tikz/fill=white!50, /tikz/minimum size=15pt}}
\begin{feynman}
\vertex (l) {};
\vertex[below = 16pt of l, dot, minimum size=3pt, label = {below: {\footnotesize $1$}}] (x1) {};
\vertex[above = 16pt of l, dot, minimum size=3pt, label = {above: {\footnotesize $2$}}] (x2) {};
\vertex[right = 8pt of l, dot, minimum size=0pt] (v12) {};
\vertex[right = 30pt of v12, blob] (b12) {};
\vertex[right = 16pt of b12, dot, minimum size=0pt] (w12) {};
\vertex[above = 1pt of w12, label = {above: {\scriptsize \hspace{-10pt} $z_{j_1}$}}] (w12u) {};
\vertex[below = 12pt of w12] (w12d) {};
\vertex[left = 10pt of w12d] (w12dl) {};
\tikzfeynmanset{every blob = {/tikz/fill=black!50, /tikz/minimum size=15pt}}
\vertex[right = 3pt of w12, blob , minimum size = 7pt] (b) {{\scriptsize $\delta$}};
\vertex[right = 3pt of b, dot, minimum size=0pt] (w34) {};
\tikzfeynmanset{every blob = {/tikz/fill=white!50, /tikz/minimum size=15pt}}
%This control the distance to the right white blob
\vertex[right = 35pt of w34, blob] (b34) {};
\vertex[above = 5pt of w34, label = {above: {\scriptsize \hspace{30pt} $\textcolor{color1}{3}$}}] (w34u) {};
\vertex[below = 12pt of w34, label = {above: {\scriptsize \hspace{30pt} $\textcolor{color1}{4}$}}] (w34d) {};
\vertex[right = 10pt of w34d] (w34dr) {};
\vertex[right = 30pt of b34, dot, minimum size=0pt] (v34) {};
\vertex[right = 30pt of b34, dot, minimum size=0pt] (vv34) {};
\vertex[right = 8pt of v34] (r) {};
\vertex[above = 16pt of r, color1, dot, minimum size=3pt, label = {above: {\footnotesize $\textcolor{color1}{3}$}}] (x3) {};
\vertex[below = 16pt of r, color1, dot, minimum size=3pt, label = {below: {\footnotesize $\textcolor{color1}{4}$}}] (x4) {};
\diagram*{
	(x1) -- (v12) -- (x2),
	(v12) -- [photon, edge label = {\scriptsize \;$\widehat{\Delta G}_{j_1}$}, inner sep = 4pt] (b12) -- [quarter left] (w12) -- [quarter left] (b12),
	(b34) -- [color1, ghost, quarter left] (w34) -- [color1, ghost, quarter left] (b34) -- [color1doubghost, edge label = {\scriptsize \,$\sigma'_{j_{2}}\sigma'_{j_{2}}$}, inner sep = 4pt] (v34), 
	(x3) -- [color1, ghost] (v34) -- [color1, ghost] (x4)
};
\end{feynman}
\end{tikzpicture}
\end{eqnarray}
or, in algebraic form
\begin{eqnarray}
\frac{1}{n_{\ell}}\delta D^{(\ell)}_{12\textcolor{color1}{34}} &=& \frac{1}{n_{\ell-1}}\sum_{\beta_{1},\beta_{2} \in \{1,2,3,4\}}(\chi_{||}^{(\ell+1)})_{12,\beta_{1}\beta_{2}}(\chi_{\perp}^{(\ell+1)})_{\textcolor{color1}{34}}\delta D^{(\ell)}_{\beta_{1}\beta_{2}\textcolor{color1}{34}}
\end{eqnarray}
In this manner, one can see that the criticality conditions studied in the main text, namely $(\chi_{||}^{(\ell+1)})_{12,\beta_{1}\beta_{2}} = \delta_{1\beta_{1}}\delta_{2\beta_{2}}$, $(\chi_{\perp}^{(\ell+1)})_{\textcolor{color1}{34}} = 1$, immediately imply the stability of $D$. 

\item Analogously, one can study the stability of the NTK variance when the network has been set to criticality. In such a case, the variation of the tensor $A$ at layer $\ell + 1$ readily follows from its corresponding recursive relation. In diagrams: 
\begin{eqnarray}
\begin{tikzpicture}[baseline=(d.base)]
\begin{feynman}
\vertex (d) {};
\vertex[left = 4pt of d, dot, minimum size=0pt] (v12) {};
\vertex[left = 16pt of v12, dot, minimum size=0pt] (l) {};
\vertex[below = 16pt of l, color1, dot, minimum size=3pt, label = {below: {\footnotesize $\textcolor{color1}{1}$}}] (x1) {};
\vertex[above = 16pt of l, color1, dot, minimum size=3pt, label = {above: {\footnotesize $\textcolor{color1}{2}$}}] (x2) {};
\vertex[right = 4pt of d, dot, minimum size=0pt] (v34) {};
\vertex[right = 16pt of v34, dot, minimum size=0pt] (r) {};
\vertex[above = 16pt of r, color2, dot, minimum size=3pt, label = {above: {\footnotesize $\textcolor{color2}{3}$}}] (x3) {};
\vertex[below = 16pt of r, color2, dot, minimum size=3pt, label = {below: {\footnotesize $\textcolor{color2}{4}$}}] (x4) {};
\diagram*{
	(x1) -- [color1, ghost] (v12) -- [color1, ghost] (x2), 
	 (x3) -- [color2, ghost] (v34) -- [color2, ghost] (x4)
};
\tikzfeynmanset{every blob = {/tikz/fill=black!50, /tikz/minimum size=15pt}}
\vertex[right = 0pt of d, blob, minimum size = 8pt] (dd) {{\scriptsize $\delta$}};
\end{feynman}
\end{tikzpicture}
&=& \sum_{j_1, j_2}
%
\begin{tikzpicture}[baseline=(b)]
\tikzfeynmanset{every blob = {/tikz/fill=white!50, /tikz/minimum size=15pt}}
\begin{feynman}
\vertex (l) {};
\vertex[below = 16pt of l, dot, color1, minimum size=3pt, label = {below: {\footnotesize $\textcolor{color1}{1}$}}] (x1) {};
\vertex[above = 16pt of l, dot, color1, minimum size=3pt, label = {above: {\footnotesize $\textcolor{color1}{2}$}}] (x2) {};
\vertex[right = 8pt of l, dot, minimum size=0pt] (v12) {};
\vertex[right = 35pt of v12, blob] (b12) {};
\vertex[right = 35pt of b12, dot, minimum size=0pt] (w12) {};
\vertex[above = 5pt of w12, label = {above: {\scriptsize \hspace{-35pt} $ \textcolor{color1}{2}$}}] (w12u) {};
\vertex[below = 12pt of w12, label = {above: {\scriptsize \hspace{-35pt} $\textcolor{color1}{1}$}}] (w12d) {};
\vertex[left = 10pt of w12d] (w12dl) {};
\tikzfeynmanset{every blob = {/tikz/fill=black!50, /tikz/minimum size=15pt}}
\vertex[right = 3pt of w12, blob , minimum size = 7pt] (b) {{\scriptsize $\delta$}};
\vertex[right = 3pt of b, dot, minimum size=0pt] (w34) {};
\tikzfeynmanset{every blob = {/tikz/fill=white!50, /tikz/minimum size=15pt}}
%This control the distance to the right white blob
\vertex[right = 35pt of w34, blob] (b34) {};
\vertex[above = 5pt of w34, label = {above: {\scriptsize \hspace{30pt} $\textcolor{color2}{3}$}}] (w34u) {};
\vertex[below = 12pt of w34, label = {above: {\scriptsize \hspace{30pt} $\textcolor{color2}{4}$}}] (w34d) {};
\vertex[right = 10pt of w34d] (w34dr) {};
\vertex[right = 30pt of b34, dot, minimum size=0pt] (v34) {};
\vertex[right = 30pt of b34, dot, minimum size=0pt] (vv34) {};
\vertex[right = 8pt of v34] (r) {};
\vertex[above = 16pt of r, color2, dot, minimum size=3pt, label = {above: {\footnotesize $\textcolor{color2}{3}$}}] (x3) {};
\vertex[below = 16pt of r, color2, dot, minimum size=3pt, label = {below: {\footnotesize $\textcolor{color2}{4}$}}] (x4) {};
\diagram*{
	(x1) -- [color1, ghost] (v12) -- [color1, ghost] (x2),
	(v12) -- [color1doubghost, edge label = {\scriptsize \;$\sigma'_{j_1}\sigma'_{j_1}$}, inner sep = 4pt] (b12) -- [color1, ghost, quarter left] (w12) -- [color1, ghost, quarter left] (b12),
	(b34) -- [color2, ghost, quarter left] (w34) -- [color2, ghost, quarter left] (b34) -- [color2doubghost, edge label = {\scriptsize \,$\sigma'_{j_{2}}\sigma'_{j_{2}}$}, inner sep = 4pt] (v34), 
	(x3) -- [color2, ghost] (v34) -- [color2, ghost] (x4)
};
\end{feynman}
\end{tikzpicture}
\end{eqnarray}
or, equivalently
\begin{eqnarray}
\frac{1}{n_{\ell}}\delta A^{(\ell)}_{\textcolor{color1}{12}\textcolor{color2}{34}} &=& \frac{1}{n_{\ell-1}}(\chi_{\perp}^{(\ell+1)})_{\textcolor{color1}{12}}(\chi_{\perp}^{(\ell+1)})_{\textcolor{color2}{34}}\delta D^{(\ell)}_{\textcolor{color1}{12}\textcolor{color2}{34}}
\end{eqnarray}
Hence, we see that the criticality condition $(\chi_{\perp}^{\ell+1})_{\alpha\beta} = 1$, directly demands the stability of the tensor $A$.


\item In similar fashion, the variation of the tensor $B$ under criticality conditions, reads
\begin{eqnarray}
\begin{tikzpicture}[baseline=(d.base)]
\begin{feynman}
\vertex (d) {};
\vertex[left = 4pt of d, dot, minimum size=0pt] (v12) {};
\vertex[left = 16pt of v12, dot, minimum size=0pt] (l) {};
\vertex[below = 16pt of l, color1, dot, minimum size=3pt, label = {below: {\footnotesize $\textcolor{color1}{1}$}}] (x1) {};
\vertex[above = 16pt of l, color2, dot, minimum size=3pt, label = {above: {\footnotesize $\textcolor{color2}{3}$}}] (x2) {};
\vertex[right = 4pt of d, dot, minimum size=0pt] (v34) {};
\vertex[right = 16pt of v34, dot, minimum size=0pt] (r) {};
\vertex[above = 16pt of r, color1, dot, minimum size=3pt, label = {above: {\footnotesize $\textcolor{color1}{2}$}}] (x3) {};
\vertex[below = 16pt of r, color2, dot, minimum size=3pt, label = {below: {\footnotesize $\textcolor{color2}{4}$}}] (x4) {};
\diagram*{
	(x1) -- [color1, ghost] (v12) -- [color2, ghost] (x2), 
	 (x3) -- [color1, ghost] (v34) -- [color2, ghost] (x4)
};
\tikzfeynmanset{every blob = {/tikz/fill=black!50, /tikz/minimum size=15pt}}
\vertex[right = 0pt of d, blob, minimum size = 8pt] (dd) {{\scriptsize $\delta$}};
\end{feynman}
\end{tikzpicture}
&=& \sum_{j_1, j_2}
%
\begin{tikzpicture}[baseline=(b)]
\tikzfeynmanset{every blob = {/tikz/fill=white!50, /tikz/minimum size=15pt}}
\begin{feynman}
\vertex (l) {};
\vertex[below = 16pt of l, dot, color1, minimum size=3pt, label = {below: {\footnotesize $\textcolor{color1}{1}$}}] (x1) {};
\vertex[above = 16pt of l, dot, color2, minimum size=3pt, label = {above: {\footnotesize $\textcolor{color2}{3}$}}] (x2) {};
\vertex[right = 8pt of l, dot, minimum size=0pt] (v12) {};
\vertex[right = 35pt of v12, blob] (b12) {};
\vertex[right = 35pt of b12, dot, minimum size=0pt] (w12) {};
\vertex[above = 5pt of w12, label = {above: {\scriptsize \hspace{-35pt} $ \textcolor{color2}{3}$}}] (w12u) {};
\vertex[below = 12pt of w12, label = {above: {\scriptsize \hspace{-35pt} $\textcolor{color1}{1}$}}] (w12d) {};
\vertex[left = 10pt of w12d] (w12dl) {};
\tikzfeynmanset{every blob = {/tikz/fill=black!50, /tikz/minimum size=15pt}}
\vertex[right = 3pt of w12, blob , minimum size = 7pt] (b) {{\scriptsize $\delta$}};
\vertex[right = 3pt of b, dot, minimum size=0pt] (w34) {};
\tikzfeynmanset{every blob = {/tikz/fill=white!50, /tikz/minimum size=15pt}}
%This control the distance to the right white blob
\vertex[right = 35pt of w34, blob] (b34) {};
\vertex[above = 5pt of w34, label = {above: {\scriptsize \hspace{30pt} $\textcolor{color1}{2}$}}] (w34u) {};
\vertex[below = 12pt of w34, label = {above: {\scriptsize \hspace{30pt} $\textcolor{color2}{4}$}}] (w34d) {};
\vertex[right = 10pt of w34d] (w34dr) {};
\vertex[right = 30pt of b34, dot, minimum size=0pt] (v34) {};
\vertex[right = 30pt of b34, dot, minimum size=0pt] (vv34) {};
\vertex[right = 8pt of v34] (r) {};
\vertex[above = 16pt of r, color1, dot, minimum size=3pt, label = {above: {\footnotesize $\textcolor{color1}{2}$}}] (x3) {};
\vertex[below = 16pt of r, color2, dot, minimum size=3pt, label = {below: {\footnotesize $\textcolor{color2}{4}$}}] (x4) {};
\diagram*{
	(x1) -- [color1, ghost] (v12) -- [color2, ghost] (x2),
	(v12) -- [color2color1ghost, edge label = {\scriptsize \;$\sigma'_{j_1}\sigma'_{j_1}$}, inner sep = 4pt] (b12) -- [color2, ghost, quarter left] (w12) -- [color1, ghost, quarter left] (b12),
	(b34) -- [color2, ghost, quarter left] (w34) -- [color1, ghost, quarter left] (b34) -- [color1color2ghost, edge label = {\scriptsize \,$\sigma'_{j_{2}}\sigma'_{j_{2}}$}, inner sep = 4pt] (v34), 
	(x3) -- [color1, ghost] (v34) -- [color2, ghost] (x4)
};
\end{feynman}
\end{tikzpicture}
\end{eqnarray}
or, explicitly
\begin{eqnarray}
\frac{1}{n_{\ell}}\delta B^{(\ell)}_{\textcolor{color1}{1}\textcolor{color2}{3}\textcolor{color1}{2}\textcolor{color2}{4}} &=& \frac{1}{n_{\ell-1}}(\chi_{\perp}^{(\ell+1)})_{\textcolor{color1}{1}\textcolor{color2}{3}}(\chi_{\perp}^{(\ell+1)})_{\textcolor{color1}{2}\textcolor{color2}{4}}\delta B^{(\ell)}_{\textcolor{color1}{1}\textcolor{color2}{3}\textcolor{color1}{2}\textcolor{color2}{4}}
\end{eqnarray}
Once more, one concludes that criticality, $(\chi_{\perp}^{\ell+1})_{\alpha\beta} = 1$, sets the stability of $B$.

% \begin{tikzpicture}[baseline=(b)]
% \tikzfeynmanset{every blob = {/tikz/fill=white!50, /tikz/minimum size=15pt}}
% \begin{feynman}
% \vertex (l) {};
% \vertex[below = 16pt of l, color1, dot, minimum size=3pt, label = {below: {\footnotesize $\textcolor{color1}{1}$}}] (x1) {};
% \vertex[above = 16pt of l, color1, dot, minimum size=3pt, label = {above: {\footnotesize $\textcolor{color1}{2}$}}] (x2) {};
% \vertex[right = 8pt of l, dot, minimum size=0pt] (v12) {};
% \vertex[right = 35pt of v12, blob] (b12) {};
% %This controls the distance of the first white blob to the black blob
% \vertex[right = 35pt of b12, dot, minimum size=0pt] (w12) {};
% \vertex[above = 5pt of w12, label = {above: {\scriptsize \hspace{-35pt} $\textcolor{color1}{2}$}}] (w12u) {};
% \vertex[below = 12pt of w12, label = {above: {\scriptsize \hspace{-35pt} $\textcolor{color1}{1}$}}] (w12d) {};
% \vertex[left = 10pt of w12d] (w12dl) {};
% \tikzfeynmanset{every blob = {/tikz/fill=gray!50, /tikz/minimum size=15pt}}
% \vertex[right = 3pt of w12, blob , minimum size = 6pt] (b) {};
% \vertex[right = 3pt of b, dot, minimum size=0pt] (w34) {};
% \tikzfeynmanset{every blob = {/tikz/fill=white!50, /tikz/minimum size=15pt}}
% %This controls the distance of the right white blob to the black blob
% \vertex[right = 35pt of w34, blob] (b34) {};
% \vertex[above = 5pt of w34, label = {above: {\scriptsize \hspace{30pt} $\textcolor{color2}{3}$}}] (w34u) {};
% \vertex[below = 12pt of w34, label = {above: {\scriptsize \hspace{30pt} $\textcolor{color2}{4}$}}] (w34d) {};
% \vertex[right = 10pt of w34d] (w34dr) {};
% \vertex[right = 35pt of b34, dot, minimum size=0pt] (v34) {};
% \vertex[right = 8pt of v34] (r) {};
% \vertex[above = 16pt of r, color2, dot, minimum size=3pt, label = {above: {\footnotesize $\textcolor{color2}{3}$}}] (x3) {};
% \vertex[below = 16pt of r, color2, dot, minimum size=3pt, label = {below: {\footnotesize $\textcolor{color2}{4}$}}] (x4) {};
% \diagram*{
% 	(x1) -- [color1, ghost] (v12) -- [color1, ghost] (x2),
% 	(v12) -- [color1doubghost, edge label = {\scriptsize \;$\sigma'_{j_1}\sigma'_{j_1}$}, inner sep = 4pt] (b12) -- [color1, ghost, quarter left] (w12) -- [color1, ghost, quarter left] (b12),
% 	(b34) -- [color2, ghost, quarter left] (w34) -- [color2, ghost, quarter left] (b34) -- [color2doubghost, edge label = {\scriptsize \,$\sigma'_{j_{2}}\sigma'_{j_{2}}$}, inner sep = 4pt] (v34),
% 	(x3) -- [color2, ghost] (v34) -- [color2, ghost] (x4)
% };
% \end{feynman}
% \end{tikzpicture}


\item One can now easily generalize our previous analysis of stability to higher-order tensors. Following our bootstrap strategy, the stability of the tensor $D^{(\ell+1)}_{1234\ldots \textcolor{color1}{(m-1)}\textcolor{color1}{m}}$ in \eqref{higher-order-tensors}, can easily be deduced from the diagrammatic relation
\begin{eqnarray}
    \begin{tikzpicture}[baseline=(d.base)]
        \begin{feynman}
        \vertex (d) {};
        \vertex[left = 4pt of d, dot, minimum size = 0pt] (v12) {};
        \vertex[left = 16pt of v12] (l) {};
        \vertex[below = 8pt of l, dot, minimum size = 3pt, label = {above: {\footnotesize $3$}}] (x1) {};
        \vertex[above = 8pt of l, dot, minimum size = 3pt, label = {above: {\footnotesize $4$}}] (x2) {};
        \vertex[below = 4pt of d, dot, minimum size = 0pt] (v34) {};
        \vertex[below = 16pt of v34] (u) {};
        \vertex[left = 8pt of u, dot, minimum size = 3pt, label = {above: {\footnotesize $2$}}] (x3) {};
        \vertex[right = 8pt of u, dot, minimum size = 3pt, label = {above: {\footnotesize $1$}}] (x4) {};
        \vertex[right = 4pt of d, dot, minimum size = 0pt] (v56) {};
        \vertex[right = 16pt of v56, dot, minimum size = 0pt] (r) {};
        \vertex[above = 8pt of r, dot, color1, minimum size = 3pt, label = {above: {\footnotesize \hspace{7pt} $\textcolor{color1}{(m-1)}$}}] (x5) {};
        \vertex[below = 8pt of r, dot, color1, minimum size = 3pt, label = {below: {\footnotesize \hspace{7pt} $\textcolor{color1}{m}$}}] (x6) {};
        \diagram*{
            (x1) -- (v12) -- (x2),
            (x3) -- (v34) -- (x4),
            (x5) -- [color1, ghost] (v56) -- [color1, ghost] (x6)
        };
        \tikzfeynmanset{every blob = {/tikz/fill=black!50, /tikz/minimum size=15pt}}
        \vertex[right = 0pt of d, blob, minimum size = 8pt] (dd) {{\scriptsize $\delta$}};
        \vertex[above = 15pt of d, dot, minimum size = 1pt] (d2) {};
        \vertex[left = 8pt of d2] (d2l) {};
        \vertex[below = 2pt of d2l, dot, minimum size = 1pt] (d1) {};
        \vertex[right = 8pt of d2] (d2r) {};
        \vertex[below = 2pt of d2r, dot, minimum size = 1pt] (d3) {};
        \end{feynman}
        \end{tikzpicture}
        %
        \;\;\,=\; \sum_{j_1, \dots, j_k}
        %
        \begin{tikzpicture}[baseline=(d.base)]
        \begin{feynman}
        \vertex (d) {};
        \vertex[left = 4pt of d, dot, minimum size = 0pt] (w12) {};
        \vertex[above = 1pt of w12, label = {above: {\scriptsize \hspace{-10pt} $z_{j_2}$}}] (w12d) {};
        \tikzfeynmanset{every blob = {/tikz/fill=white!50, /tikz/minimum size=15pt}}
        \vertex[left = 25pt of w12, blob] (b12) {};
        \vertex[left = 25pt of b12, dot, minimum size = 0pt] (v12) {};
        \vertex[left = 8pt of v12] (l) {};
        \vertex[below = 16pt of l, dot, minimum size = 3pt, label = {below: {\footnotesize $3$}}] (x1) {};
        \vertex[above = 16pt of l, dot, minimum size = 3pt, label = {above: {\footnotesize $4$}}] (x2) {};
        \vertex[below = 4pt of d, dot, minimum size = 0pt] (w34) {};
        \vertex[right = 13pt of w34, label = {above: {\scriptsize $ $}}, inner sep = -1pt] (w34r) {};
        \tikzfeynmanset{every blob = {/tikz/fill=white!50, /tikz/minimum size=15pt}}
        \vertex[below = 25pt of w34, blob] (b34) {};
        \vertex[below = 25pt of b34, dot, minimum size = 0pt] (v34) {};
        \vertex[below = 8pt of v34] (u) {};
        \vertex[left = 16pt of u, dot, minimum size = 3pt, label = {above: {\footnotesize $2$}}] (x3) {};
        \vertex[right = 16pt of u, dot, minimum size = 3pt, label = {above: {\footnotesize $1$}}] (x4) {};
        \vertex[right = 4pt of d, dot, minimum size = 0pt] (w56) {};
        \vertex[below = 1pt of w56, label = {below: {\scriptsize \hspace{18pt} $z_{j_1}$}}] (w56d) {};
        \tikzfeynmanset{every blob = {/tikz/fill=white!50, /tikz/minimum size=15pt}}
        \vertex[right = 25pt of w56, blob] (b56) {};
        \vertex[right = 25pt of b56, dot, minimum size = 0pt] (v56) {};
        \vertex[right = 8 pt of v56, dot, minimum size = 0pt] (r) {};
        \vertex[above = 16pt of r, dot, color1, minimum size = 3pt, label = {above: {\footnotesize $\textcolor{color1}{(m-1)}$}}] (x5) {};
        \vertex[below = 16pt of r, dot, color1, minimum size = 3pt, label = {below: {\footnotesize $\textcolor{color1}{m}$}}] (x6) {};
        \diagram*{
            (x1) -- (v12) -- (x2),
            (x3) -- (v34) -- (x4),
            (x5) -- [color1, ghost] (v56) -- [color1, ghost] (x6),
            (v12) -- [photon, edge label' = {\scriptsize \;$\widehat{\Delta G}_{j_2}$}, inner sep = 4pt] (b12) -- [quarter left] (w12) -- [quarter left] (b12),
            (v34) -- [photon, edge label = {\scriptsize $\widehat{\Delta G}_{j_1}$}, inner sep = 4pt] (b34) -- [quarter left] (w34) -- [quarter left] (b34),
            (v56) -- [color1doubghost, edge label = {\scriptsize \,$\sigma'_{j_k}\sigma'_{j_k}$}, inner sep = 4pt] (b56) -- [color1, ghost, quarter left] (w56) -- [color1, ghost, quarter left] (b56)
        };
        \tikzfeynmanset{every blob = {/tikz/fill=black!50, /tikz/minimum size=15pt}}
        \vertex[right = 0pt of d, blob, minimum size = 8pt] (dd) {{\scriptsize $\delta$}};
        \vertex[above = 25pt of d, dot, minimum size = 1pt] (d2) {};
        \vertex[left = 10pt of d2] (d2l) {};
        \vertex[below = 3pt of d2l, dot, minimum size = 1pt] (d1) {};
        \vertex[right = 10pt of d2] (d2r) {};
        \vertex[below = 3pt of d2r, dot, minimum size = 1pt] (d3) {};
        \end{feynman}
        \end{tikzpicture}
\end{eqnarray}
which algebraically reads
\begin{eqnarray}\label{dequationcriticality}
    \delta D^{(\ell+1)}_{1234\ldots \textcolor{color1}{(m-1)}\textcolor{color1}{m}} \! \! \! \! \! \! &=& \! \! \! \! \! \! \sum_{\beta_{k}}(\chi^{(\ell+1)}_{\parallel})_{12,\beta_{1}\beta_{2}}(\chi^{(\ell+1)}_{\parallel})_{34,\beta_{3}\beta_{4}}\ldots (\chi^{(\ell+1)}_{\perp})_{\textcolor{color1}{(m-1)}\textcolor{color1}{m}}\delta D^{(\ell)}_{\beta_{1}\beta_{2}\beta_{3}\beta_{4}\ldots \textcolor{color1}{(m-1)}\textcolor{color1}{m}}\nonumber\\ 
\end{eqnarray}
Remarkably, one sees that the criticality conditions for $\chi_{\parallel}$ in~\eqref{criticalityconstraintsequiv} and $\chi_{\perp}$ below~\eqref{eq:8} are again sufficient to fix the stability of the tensor $D$.

\item The higher-order generalization of the tensors $A$ and $B$ are obtained from the following expectation value involving two NTks:
 \begin{align}\label{higher-order-tensors_AB}
     &\mathbb{E}^{c}_{\theta}[z^{(\ell+1)}_{i_{1},1},z^{(\ell+1)}_{i_{2},2},\widehat{\Delta \Theta}^{(\ell+1)}_{i_{3}i_{4},\textcolor{color1}{3}\textcolor{color1}{4}} ,\ldots, \widehat{\Delta \Theta}^{(\ell+1)}_{i_{m-1}i_{m},\textcolor{color2}{(m-1)}\textcolor{color2}{m}}] \nonumber\\ 
      &{=} \frac{1}{n_{\ell}^{\frac{m}{2}-1}}\!\bigg[\delta_{i_{1}i_{2}}\delta_{i_{3}i_{4}}\ldots \delta_{i_{m-1}i_{m}}A^{(\ell+1)}_{12\textcolor{color1}{34}\ldots \textcolor{color2}{(m-1)m}} {+} \delta_{i_{1}i_{2}}\delta_{i_{3}i_{(m-1)}}\ldots \delta_{i_{4}i_{m}}B^{(\ell+1)}_{12\textcolor{color1}{3}\textcolor{color2}{(m-1)}\ldots \textcolor{color1}{4}\textcolor{color2}{m}} {+} \ldots\!\bigg]
\end{align}
The stability of the higher-order tensors $A$ and $B$ can easily be studied using our diagrammatic approach. The variation of the tensor $A$ takes the form
\begin{eqnarray}
    \begin{tikzpicture}[baseline=(d.base)]
        \begin{feynman}
        \vertex (d) {};
        \vertex[left = 4pt of d, dot, minimum size = 0pt] (v12) {};
        \vertex[left = 16pt of v12] (l) {};
        \vertex[below = 8pt of l, dot, color1, minimum size = 3pt, label = {above: {\footnotesize $\textcolor{color1}{3}$}}] (x1) {};
        \vertex[above = 8pt of l, dot, color1, minimum size = 3pt, label = {above: {\footnotesize $\textcolor{color1}{4}$}}] (x2) {};
        \vertex[below = 4pt of d, dot, minimum size = 0pt] (v34) {};
        \vertex[below = 16pt of v34] (u) {};
        \vertex[left = 8pt of u, dot, minimum size = 3pt, label = {above: {\footnotesize $2$}}] (x3) {};
        \vertex[right = 8pt of u, dot, minimum size = 3pt, label = {above: {\footnotesize $1$}}] (x4) {};
        \vertex[right = 4pt of d, dot, minimum size = 0pt] (v56) {};
        \vertex[right = 16pt of v56, dot, minimum size = 0pt] (r) {};
        \vertex[above = 8pt of r, dot, color2, minimum size = 3pt, label = {above: {\footnotesize \hspace{7pt} $\textcolor{color2}{(m-1)}$}}] (x5) {};
        \vertex[below = 8pt of r, dot, color2, minimum size = 3pt, label = {below: {\footnotesize \hspace{7pt} $\textcolor{color2}{m}$}}] (x6) {};
        \diagram*{
            (x1) -- [color1, ghost] (v12) -- [color1, ghost] (x2),
            (x3) -- (v34) -- (x4),
            (x5) -- [color2, ghost] (v56) -- [color2, ghost] (x6)
        };
        \tikzfeynmanset{every blob = {/tikz/fill=black!50, /tikz/minimum size=15pt}}
        \vertex[right = 0pt of d, blob, minimum size = 8pt] (dd) {{\scriptsize $\delta$}};
        \vertex[above = 15pt of d, dot, minimum size = 1pt] (d2) {};
        \vertex[left = 8pt of d2] (d2l) {};
        \vertex[below = 2pt of d2l, dot, minimum size = 1pt] (d1) {};
        \vertex[right = 8pt of d2] (d2r) {};
        \vertex[below = 2pt of d2r, dot, minimum size = 1pt] (d3) {};
        \end{feynman}
        \end{tikzpicture}
        %
        \;\;\,=\; \sum_{j_1, \dots, j_k}
        %
        \begin{tikzpicture}[baseline=(d.base)]
        \begin{feynman}
        \vertex (d) {};
        \vertex[left = 4pt of d, dot, minimum size = 0pt] (w12) {};
        \vertex[above = 1pt of w12, label = {above: {\scriptsize \hspace{-10pt} $ $}}] (w12d) {};
        \tikzfeynmanset{every blob = {/tikz/fill=white!50, /tikz/minimum size=15pt}}
        \vertex[left = 25pt of w12, blob] (b12) {};
        \vertex[left = 25pt of b12, dot, minimum size = 0pt] (v12) {};
        \vertex[left = 8pt of v12] (l) {};
        \vertex[below = 16pt of l, dot, color1, minimum size = 3pt, label = {below: {\footnotesize $\textcolor{color1}{3}$}}] (x1) {};
        \vertex[above = 16pt of l, dot, color1, minimum size = 3pt, label = {above: {\footnotesize $\textcolor{color1}{4}$}}] (x2) {};
        \vertex[below = 4pt of d, dot, minimum size = 0pt] (w34) {};
        \vertex[right = 13pt of w34, label = {above: {\scriptsize $ $}}, inner sep = -1pt] (w34r) {};
        \tikzfeynmanset{every blob = {/tikz/fill=white!50, /tikz/minimum size=15pt}}
        \vertex[below = 25pt of w34, blob] (b34) {};
        \vertex[below = 25pt of b34, dot, minimum size = 0pt] (v34) {};
        \vertex[below = 8pt of v34] (u) {};
        \vertex[left = 16pt of u, dot, minimum size = 3pt, label = {above: {\footnotesize $2$}}] (x3) {};
        \vertex[right = 16pt of u, dot, minimum size = 3pt, label = {above: {\footnotesize $1$}}] (x4) {};
        \vertex[right = 4pt of d, dot, minimum size = 0pt] (w56) {};
        \vertex[below = 1pt of w56, label = {below: {\scriptsize \hspace{18pt} $z_{j_1} $}}] (w56d) {};
        \tikzfeynmanset{every blob = {/tikz/fill=white!50, /tikz/minimum size=15pt}}
        \vertex[right = 25pt of w56, blob] (b56) {};
        \vertex[right = 25pt of b56, dot, minimum size = 0pt] (v56) {};
        \vertex[right = 8 pt of v56, dot, minimum size = 0pt] (r) {};
        \vertex[above = 16pt of r, dot, color2, minimum size = 3pt, label = {above: {\footnotesize $\textcolor{color2}{(m-1)}$}}] (x5) {};
        \vertex[below = 16pt of r, dot, color2, minimum size = 3pt, label = {below: {\footnotesize $\textcolor{color2}{m}$}}] (x6) {};
        \diagram*{
            (x1) -- [color1, ghost] (v12) -- [color1, ghost] (x2),
            (x3) -- (v34) -- (x4),
            (x5) -- [color2, ghost] (v56) -- [color2, ghost] (x6),
            (v12) -- [color1doubghost, edge label' = {\scriptsize \;$\sigma'_{j_2}\sigma'_{j_2}$}, inner sep = 4pt] (b12) -- [color1, ghost, quarter left] (w12) -- [color1, ghost, quarter left] (b12),
            (v34) -- [photon, edge label = {\scriptsize $\widehat{\Delta G}_{j_1}$}, inner sep = 4pt] (b34) -- [quarter left] (w34) -- [quarter left] (b34),
            (v56) -- [color2doubghost, edge label = {\scriptsize \,$\sigma'_{j_k}\sigma'_{j_k}$}, inner sep = 4pt] (b56) -- [color2, ghost, quarter left] (w56) -- [color2, ghost, quarter left] (b56)
        };
        \tikzfeynmanset{every blob = {/tikz/fill=black!50, /tikz/minimum size=15pt}}
        \vertex[right = 0pt of d, blob, minimum size = 8pt] (dd) {{\scriptsize $\delta$}};
        \vertex[above = 25pt of d, dot, minimum size = 1pt] (d2) {};
        \vertex[left = 10pt of d2] (d2l) {};
        \vertex[below = 3pt of d2l, dot, minimum size = 1pt] (d1) {};
        \vertex[right = 10pt of d2] (d2r) {};
        \vertex[below = 3pt of d2r, dot, minimum size = 1pt] (d3) {};
        \end{feynman}
        \end{tikzpicture}
\end{eqnarray}
or, in analytical form
\begin{eqnarray}
    \delta A^{(\ell+1)}_{12\textcolor{color1}{34}\ldots \textcolor{color2}{(m-1)m}} &=& \sum_{\beta_{k}}(\chi_{||}^{(\ell+1)})_{12,\beta_{1}\beta_{2}}(\chi_{\perp}^{(\ell+1)})_{\textcolor{color1}{34}}\ldots (\chi_{\perp}^{(\ell+1)})_{\textcolor{color2}{(m-1)m}}\delta A^{(\ell)}_{\beta_{1}\beta_{2}\textcolor{color1}{34}\ldots \textcolor{color2}{(m-1)m}} \nonumber\\
\end{eqnarray}
which shows $A$ stays stable once the criticality conditions studied previously are satisfied. 

Similarly, the variation of the tensor $B$ reads 
\begin{eqnarray}
    \begin{tikzpicture}[baseline=(d.base)]
        \begin{feynman}
        \vertex (d) {};
        \vertex[left = 4pt of d, dot, minimum size = 0pt] (v12) {};
        \vertex[left = 16pt of v12] (l) {};
        \vertex[below = 8pt of l, dot, color1, minimum size = 3pt, label = {above: {\footnotesize $\textcolor{color1}{3}$}}] (x1) {};
        \vertex[above = 8pt of l, dot, color2, minimum size = 3pt, label = {above: {\footnotesize $\textcolor{color2}{(m-1)}$}}] (x2) {};
        \vertex[below = 4pt of d, dot, minimum size = 0pt] (v34) {};
        \vertex[below = 16pt of v34] (u) {};
        \vertex[left = 8pt of u, dot, minimum size = 3pt, label = {above: {\footnotesize $2$}}] (x3) {};
        \vertex[right = 8pt of u, dot, minimum size = 3pt, label = {above: {\footnotesize $1$}}] (x4) {};
        \vertex[right = 4pt of d, dot, minimum size = 0pt] (v56) {};
        \vertex[right = 16pt of v56, dot, minimum size = 0pt] (r) {};
        \vertex[above = 8pt of r, dot, color1, minimum size = 3pt, label = {above: {\footnotesize \hspace{7pt} $\textcolor{color1}{4}$}}] (x5) {};
        \vertex[below = 8pt of r, dot, color2, minimum size = 3pt, label = {below: {\footnotesize \hspace{7pt} $\textcolor{color2}{m}$}}] (x6) {};
        \diagram*{
            (x1) -- [color1, ghost] (v12) -- [color2, ghost] (x2),
            (x3) -- (v34) -- (x4),
            (x5) --[color1, ghost] (v56) -- [color2, ghost] (x6)
        };
        \tikzfeynmanset{every blob = {/tikz/fill=black!50, /tikz/minimum size=15pt}}
        \vertex[right = 0pt of d, blob, minimum size = 8pt] (dd) {{\scriptsize $\delta$}};
        \vertex[above = 15pt of d, dot, minimum size = 1pt] (d2) {};
        \vertex[left = 8pt of d2] (d2l) {};
        \vertex[below = 2pt of d2l, dot, minimum size = 1pt] (d1) {};
        \vertex[right = 8pt of d2] (d2r) {};
        \vertex[below = 2pt of d2r, dot, minimum size = 1pt] (d3) {};
        \end{feynman}
        \end{tikzpicture}
        %
        \;\;\,=\; \sum_{j_1, \dots, j_k}
        %
        \begin{tikzpicture}[baseline=(d.base)]
        \begin{feynman}
        \vertex (d) {};
        \vertex[left = 4pt of d, dot, minimum size = 0pt] (w12) {};
        \vertex[above = 1pt of w12, label = {above: {\scriptsize \hspace{-10pt} $ $}}] (w12d) {};
        \tikzfeynmanset{every blob = {/tikz/fill=white!50, /tikz/minimum size=15pt}}
        \vertex[left = 25pt of w12, blob] (b12) {};
        \vertex[left = 25pt of b12, dot, minimum size = 0pt] (v12) {};
        \vertex[left = 8pt of v12] (l) {};
        \vertex[below = 16pt of l, dot, color1, minimum size = 3pt, label = {below: {\footnotesize $\textcolor{color1}{3}$}}] (x1) {};
        \vertex[above = 16pt of l, dot, color2, minimum size = 3pt, label = {above: {\footnotesize $\textcolor{color2}{(m-1)}$}}] (x2) {};
        \vertex[below = 4pt of d, dot, minimum size = 0pt] (w34) {};
        \vertex[right = 13pt of w34, label = {above: {\scriptsize $ $}}, inner sep = -1pt] (w34r) {};
        \tikzfeynmanset{every blob = {/tikz/fill=white!50, /tikz/minimum size=15pt}}
        \vertex[below = 25pt of w34, blob] (b34) {};
        \vertex[below = 25pt of b34, dot, minimum size = 0pt] (v34) {};
        \vertex[below = 8pt of v34] (u) {};
        \vertex[left = 16pt of u, dot, minimum size = 3pt, label = {above: {\footnotesize $2$}}] (x3) {};
        \vertex[right = 16pt of u, dot, minimum size = 3pt, label = {above: {\footnotesize $1$}}] (x4) {};
        \vertex[right = 4pt of d, dot, minimum size = 0pt] (w56) {};
        \vertex[below = 1pt of w56, label = {below: {\scriptsize \hspace{18pt} $z_{j_1} $}}] (w56d) {};
        \tikzfeynmanset{every blob = {/tikz/fill=white!50, /tikz/minimum size=15pt}}
        \vertex[right = 25pt of w56, blob] (b56) {};
        \vertex[right = 25pt of b56, dot, minimum size = 0pt] (v56) {};
        \vertex[right = 8 pt of v56, dot, minimum size = 0pt] (r) {};
        \vertex[above = 16pt of r, dot, color1, minimum size = 3pt, label = {above: {\footnotesize $\textcolor{color1}{4}$}}] (x5) {};
        \vertex[below = 16pt of r, dot, color2, minimum size = 3pt, label = {below: {\footnotesize $\textcolor{color2}{m}$}}] (x6) {};
        \diagram*{
            (x1) -- [color1, ghost] (v12) -- [color2, ghost] (x2),
            (x3) -- (v34) -- (x4),
            (x5) -- [color1, ghost] (v56) -- [color2, ghost] (x6),
            (v12) -- [color2color1ghost, edge label' = {\scriptsize \;$\sigma'_{j_2}\sigma'_{j_2}$}, inner sep = 4pt] (b12) -- [color2, ghost, quarter left] (w12) -- [color1, ghost, quarter left] (b12),
            (v34) -- [photon, edge label = {\scriptsize $\widehat{\Delta G}_{j_1}$}, inner sep = 4pt] (b34) -- [quarter left] (w34) -- [quarter left] (b34),
            (v56) -- [color1color2ghost, edge label = {\scriptsize \,$\sigma'_{j_k}\sigma'_{j_k}$}, inner sep = 4pt] (b56) -- [color2, ghost, quarter left] (w56) -- [color1, ghost, quarter left] (b56)
        };
        \tikzfeynmanset{every blob = {/tikz/fill=black!50, /tikz/minimum size=15pt}}
        \vertex[right = 0pt of d, blob, minimum size = 8pt] (dd) {{\scriptsize $\delta$}};
        \vertex[above = 25pt of d, dot, minimum size = 1pt] (d2) {};
        \vertex[left = 10pt of d2] (d2l) {};
        \vertex[below = 3pt of d2l, dot, minimum size = 1pt] (d1) {};
        \vertex[right = 10pt of d2] (d2r) {};
        \vertex[below = 3pt of d2r, dot, minimum size = 1pt] (d3) {};
        \end{feynman}
        \end{tikzpicture}
\end{eqnarray}
which translates into the formal expression
\begin{eqnarray}
        \delta B^{(\ell+1)}_{12\textcolor{color1}{3}\textcolor{color2}{(m-1)}\ldots \textcolor{color1}{4}\textcolor{color2}{m}} &=& \sum_{\beta_{k}}(\chi_{||}^{(\ell+1)})_{12,\beta_{1}\beta_{2}}(\chi_{\perp}^{(\ell+1)})_{\textcolor{color1}{3}\textcolor{color2}{(m-1)}}\ldots (\chi_{\perp}^{(\ell+1)})_{\textcolor{color1}{4}\textcolor{color2}{m}}\delta B^{(\ell)}_{\beta_{1}\beta_{2}\textcolor{color1}{3}\textcolor{color2}{(m-1)}\ldots \textcolor{color1}{4}\textcolor{color2}{m}} \nonumber\\
\end{eqnarray}
One sees again that the criticality conditions guarantee the stability of the higher-order tensor $B$.


\end{enumerate}
%%% Local Variables:
%%% mode: LaTeX
%%% TeX-master: "ntk_feynman"
%%% End:
